\section{Evaluation}
\label{sec:eval}

The evaluation asks four deliberately bounded questions.  \textbf{RQ1:} Do
independent replicas exhibit the advertised disjoint-fragment merge,
same-fragment strain, resolution, and transport behavior in controlled
end-to-end scenarios?  \textbf{RQ2:} Does the native kernel path remain
competitive with the conventional Linux block stack for sequential I/O?
\textbf{RQ3:} What portability cost does the FUSE path impose relative to a
FUSE filesystem with comparable encryption?  And \textbf{RQ4:} are the
reported throughput samples free of the process leaks, out-of-memory
terminations, and silent zero results that can invalidate a long storage
campaign?  We do not use these experiments to answer metadata-workload,
multi-client scalability, or non-loopback LAN/WAN synchronization questions.

\subsection{Functional Multi-Replica Evidence}
\label{subsec:eval-mechanisms}

The RC1 test matrix includes deterministic end-to-end scenarios for the
distributed mechanisms, summarized in \cref{tab:mechanism-evidence}.  Each has
a binary acceptance oracle over exact bytes, frontier cardinality, or an
expected authorization result; none derives latency, throughput, conflict
probability, or scale from the scenario.

\begin{table}[t]
  \centering
  \small
  \setlength{\tabcolsep}{4pt}
  \caption{Controlled mechanism scenarios in the RC1 artifact.  These are
  correctness fixtures, not workload or scalability measurements.}
  \label{tab:mechanism-evidence}
  \begin{tabular}{@{}p{0.19\linewidth}p{0.22\linewidth}p{0.50\linewidth}@{}}
    \toprule
    Mechanism & Controlled topology & Acceptance oracle \\
    \midrule
    Disjoint writes &
      two replicas, opaque local transport &
      both replicas read the exact combined fragments, expose no conflict, and
      remain stable on a second sync \\
    Same-fragment writes &
      two clients, hosted TLS service &
      both retain two byte-exact strains; a selected resolution propagates and
      collapses the remote frontier to one \\
    Late observer &
      two writers plus a fresh third client, hosted TLS service &
      the third client reconstructs exact units and both unresolved strains;
      a later fresh client observes the resolved frontier \\
    Live-peer topology &
      two containers, TLS on ephemeral loopback &
      one sync transfers each peer's distinct unit in both directions;
      wrong-key and unauthenticated access are rejected \\
    Hosted authorization &
      WriterSet clients and enforcing service &
      member frontiers are accepted; non-member and revoked-writer frontiers
      are rejected; stored-byte scans contain no injected plaintext markers \\
    \bottomrule
  \end{tabular}
\end{table}

The fixtures are implemented under \path{crates/sfs-sync/tests} and
\path{crates/sfs-saas/tests} as
\texttt{conflict\_e2e}, \texttt{net\_e2e},
\texttt{p2p\_net\_e2e}, \texttt{enforcement\_e2e}, and
\texttt{zk\_proof\_e2e}.  Together they answer RQ1 at the mechanism level:
the advertised reconciliation outcomes cross both an opaque hosted transport
and a live network peer, and an unresolved frontier remains discoverable by a
client that did not participate in the conflict.  They do not answer how often
those outcomes occur in a repository or how their cost grows with agents,
files, latency, or history.

The artifact also guards the re-chunk boundary called out by the geometry
design.  Functional fixtures grow a unit from 4-KiB to 256-KiB fragments,
require byte-exact content, fresh causal dots, preserved pinned history,
publication-safe failure, and successful reopen.  A separate physical-footprint
regression appends 512-KiB chunks to 32 MiB.  The fixture records a before/after
observation of 62 MiB of allocated blocks, or \(1.94\times\) logical size,
versus 146 MiB (\(4.58\times\)) for the superseded implementation that retained every
re-chunked generation; the guard requires less than \(2.5\times\).  This
\texttt{st\_blocks} high-water proxy includes freed but not hole-punched
intermediate generations.  It is neither a transition-latency result nor a
long-term aging experiment, and the paper does not treat it as part of the
fault-gated throughput campaign below.

\subsection{Method}
\label{subsec:eval-method}

We compare within execution surfaces, not across them.  The kernel module uses
a raw block device and is compared with ext4 on that device for plaintext and
ext4 over LUKS for encrypted I/O.  The FUSE daemon stores an \sfs{} container
on ext4 and is compared with fuse2fs for plaintext and gocryptfs for encrypted
I/O.  Thus each ratio has a baseline in the same kernel/userspace class; a
kernel-to-FUSE ratio would answer a different question.

The published campaign used kernel artifact \texttt{4cb3e28} and FUSE artifact
\texttt{4a41dbeb}.  It ran on a Linux benchmark VM backed by a Samsung 960 EVO,
using \texttt{fio}'s single-threaded \texttt{psync} engine.  Each cell is the
arithmetic mean of ten fault-free runs.  The 14 request sizes range from 4 KiB
through 2047 MiB; files are 1 GiB except for the 2047-MiB request, which uses a
2-GiB file.  The 2047-MiB endpoint stays below Linux's
\texttt{MAX\_RW\_COUNT}; literal 2- and 4-GiB single requests fail for ext4 as
well and are therefore not \sfs{} limits.

Every run starts with block discard, cache dropping, and a fresh mount.  Read
results are consequently cold-cache results.  The fixture inspects kernel
messages, OOM events, leaked daemons, and available memory after each run and
excludes a run on any such fault.  The final campaign reported zero faults, so
all ten samples entered every cell.  Source, on-filesystem, and cold-read hashes
were also compared for the plaintext, XTS, and GCM paths before accepting the
campaign.

The checked-in aggregate is
\texttt{docs/perf/perf-report-2026-07-20.html}.  For compact comparison,
Tables~\ref{tab:kernel-perf} and
\ref{tab:fuse-perf} report the geometric mean of the 14 per-size ratios
\[
  R_{\mathrm{geo}} =
  \exp\!\left(\frac{1}{14}\sum_{i=1}^{14}
  \ln\frac{\overline{x}_{\sfs{},i}}{\overline{x}_{\mathrm{base},i}}\right),
\]
where each \(\overline{x}\) is the arithmetic mean for one cell.  This gives
equal multiplicative weight to each request size.  The repository report
contains every per-size cell, but currently archives aggregates rather than
the ten raw samples and their health logs; this prevents confidence intervals
or an independent reanalysis of run-level variance.

\subsection{Kernel Path}
\label{subsec:eval-kernel}

\begin{table}[t]
  \centering
  \caption{Kernel-path throughput relative to its group baseline.  Values are
  geometric means over 14 request sizes; larger is better.}
  \label{tab:kernel-perf}
  \begin{tabular}{lllr}
    \toprule
    Content & Operation & Baseline & \sfs{}/baseline \\
    \midrule
    Plain & write & ext4-raw & \(1.22\times\) \\
    Plain & cold read & ext4-raw & \(0.98\times\) \\
    XTS & write & LUKS-ext4 & \(0.99\times\) \\
    GCM & write & LUKS-ext4 & \(0.98\times\) \\
    XTS & cold read & LUKS-ext4 & \(1.20\times\) \\
    GCM & cold read & LUKS-ext4 & \(1.11\times\) \\
    \bottomrule
  \end{tabular}
\end{table}

Table~\ref{tab:kernel-perf} answers RQ2 for this sequential, queue-depth-one
workload.  Plain writes are faster than ext4 on average, while plain reads and
encrypted writes are at parity under a \(0.98\times\) convention.  Encrypted
cold reads exceed the LUKS-ext4 baseline for both ciphers.  These are
geometric summaries, not universal dominance: individual cells can lose.  In
particular, the 2047-MiB encrypted-write endpoint is approximately
\(0.81\times\) LUKS.  The result supports the narrower claim that native
fragment versioning and authenticated metadata do not introduce an inherent
sequential-throughput penalty on this device and workload.

\subsection{FUSE Path}
\label{subsec:eval-fuse}

\begin{table}[t]
  \centering
  \caption{FUSE-path throughput relative to its group baseline, computed as in
  Table~\ref{tab:kernel-perf}.}
  \label{tab:fuse-perf}
  \begin{tabular}{lllr}
    \toprule
    Content & Operation & Baseline & \sfs{}/baseline \\
    \midrule
    Plain & write & fuse2fs & \(3.11\times\) \\
    Plain & cold read & fuse2fs & \(1.70\times\) \\
    XTS & write & gocryptfs & \(0.87\times\) \\
    GCM & write & gocryptfs & \(0.85\times\) \\
    XTS & cold read & gocryptfs & \(0.38\times\) \\
    GCM & cold read & gocryptfs & \(0.33\times\) \\
    \bottomrule
  \end{tabular}
\end{table}

The portable path has two distinct regimes.  Sequence-detecting read-ahead
lets plaintext \sfs{} outperform fuse2fs for both reads and writes, showing
that the FUSE boundary alone is not the limiting factor.  Encrypted writes are
13--15\% below gocryptfs.  Encrypted cold reads are the clear remaining
weakness: XTS reaches \(0.38\times\) and GCM \(0.33\times\) of gocryptfs in the
geometric mean.  At large requests the corresponding ratios rise only to
approximately \(0.43\times\) and \(0.38\times\).

Profiling attributes approximately 40\% of the encrypted-read time to
parallel AES-NI decryption, 19\% to container reads, and 32\% to FUSE reply,
system-call, and backpressure work, with the remainder in buffer movement.
This is a mixed end-to-end userspace ceiling, not evidence for one removable
\texttt{splice} or library bottleneck: the plaintext path reaches roughly
1.5 GB/s through the same reply implementation, and the compared gocryptfs
read path does not establish a zero-copy explanation.  Closing the gap would
therefore require a separate buffer-ownership and pipeline study.  The result
answers RQ3 honestly: FUSE is a useful portable path, but it is not the
performance path for encrypted sequential reads in this release.

\subsection{Fault Discipline and Scope}
\label{subsec:eval-validity}

A preliminary campaign was discarded in full after externally unmounted FUSE
daemons failed to exit, accumulated, and eventually caused OOM kills and
zero-valued cells.  No number from that campaign appears here.  The daemon exit
was corrected and the per-run health gate described above was added before the
reported campaign.  This experience is why RQ4 is part of the evaluation: ten
repetitions do not improve a result if failed runs silently enter the mean.

The remaining validity limits are material.  The campaign uses one VM, one
SSD, buffered single-threaded sequential I/O, and throughput rather than tail
latency.  It does not exercise random I/O, metadata-heavy application traces,
multiple queue depths, many writers, network delay, link rate,
synchronization volume, or conflict frequency.  The data therefore validate a
local data path and expose one FUSE weakness; they do not by themselves
validate the broader agentic-time workload at scale.
